\documentclass[12pt]{article}
\usepackage[T1]{fontenc}
\usepackage{amssymb}
\usepackage{amsmath}
\usepackage{color}
\usepackage{graphicx}
\usepackage{mdframed}
\usepackage{listings, xcolor}
\usepackage{booktabs}
% \usepackage[dvipsnames]{xcolor} % Non va!
\usepackage{textcomp}
\usepackage[table]{xcolor}
\newcommand{\gr}{\rowcolor[HTML]{EFEFEF}}
\lstdefinelanguage{myC}{
	keywords=[1]{break, case, continue, default, do
		, else, false, for, if, const, return, switch, true, while}, % generic keywords
	keywordstyle=[1]\color{blue}\bfseries,
	keywords=[2]{bool, int, long, float, double, byte, short, char, void, signed, unsigned}, % types
	keywordstyle=[2]\color{teal}\bfseries,
	keywordstyle=[2]\color{violet}\bfseries,
	keywords=[3]{NULL},
	keywordstyle=[3]\color{violet}\bfseries,
	identifierstyle=\color{black},
	sensitive=false,
	comment=[l]{//},
	morecomment=[s]{/*}{*/},
	commentstyle=\color{violet}\ttfamily,
	commentstyle=\normalsize\ttfamily\color{grey}, % scriptsize
	stringstyle=\color{forestgreen}\ttfamily,
	morestring=[b]',
	morestring=[b]"
}

\lstset{
	language=myC,
	backgroundcolor=\color{veryLightgray},
	extendedchars=true,
	basicstyle=\small\ttfamily,
	showstringspaces=false,
	showspaces=false,
	numbers=none,
	numberstyle=\normalsize,
	numbersep=9pt,
	tabsize=2,
	upquote=true,
	breaklines=true,
	showtabs=false,
	captionpos=b
}

\definecolor{myBlue}{rgb}{0.5,0.5,1}
\definecolor{myLightBlue}{rgb}{0.35,0.6,0.8}
\definecolor{myBlack}{rgb}{0,0,0}
\definecolor{myGreen}{rgb}{0.1,0.6,0.2}
\definecolor{myGray}{rgb}{0.5,0.5,0.5}
\definecolor{myLightgray}{rgb}{0.95,0.95,0.95}
\definecolor{verylightgray}{rgb}{.97,.97,.97}
\definecolor{myMauve}{rgb}{0.58,0,0.82}
\definecolor{forestgreen}{rgb}{0.13, 0.55, 0.13}
\definecolor{forestgreen}{rgb}{0.13, 0.55, 0.13}
\lstdefinelanguage{customc}{
    language=C,
    backgroundcolor = \color{myLightgray},
    basicstyle=\small\ttfamily\color{myBlack},
    keywordstyle=\color{myLightBlue},
    keywordstyle=[2]\color{red},
    commentstyle=\small\ttfamily\color{myGreen},
    morekeywords={RequirePackage,ProvidesPackage},
    %
    % The special highlighting works for '!if', '!endif' and '!else'
    % But it doesn't work for '#if', '#endif' and '#else'.
    alsoletter = {!},
    keywords=[2]{!if,!endif,!else},
    otherkeywords={define,include,\# }
}

\lstdefinestyle{myCustomc}{
    language = customc,
    % keywordstyle = \color{myMauve},
}

\lstset{escapechar=@,style=myCustomc}


\definecolor{grey}{rgb}{0.3,0.3,0.3}
\definecolor{lightgrey}{rgb}{0.9,0.9,0.9}

\thispagestyle{empty}
\setlength{\textwidth}{18.5cm}
\setlength{\topmargin}{-2.5cm}
\setlength{\textheight}{24.5cm}
\setlength{\oddsidemargin}{-1cm}
\setlength{\evensidemargin}{-1cm}

\begin{document}
\begin{center}
  {\LARGE Compito di Programmazione - Bioinformatica}\\
  {10 febbraio 2026 (tempo disponibile: 2 ore)}
\end{center}

\vspace*{1ex}
\begin{center}{\Large Esercizio 1} ($11$ punti)\\
  (\textbf{si consegni} il file \texttt{matrice.c})
\end{center}

Si completi il seguente programma \texttt{matrice.c}
in modo che legga da tastiera una matrice $N\times M$
di interi sicuramente positivi e quindi inizializzi un vettore $X$ di $N$ elementi,
in modo tale che
ogni elemento $X[i]$ contenga il valore $1$ se ciascun elemento della riga $i$-esima
ha almeno una cifra pari, e contenga $0$ altrimenti.

Per esempio, se $N = 3$, $M = 4$ e la matrice inserita fosse:
%
\[
\begin{bmatrix}
  12 & 32 & 46 & 27 \\
  2 & 7 & 3 & 72 \\
  5 & 9 & 77 & 11
\end{bmatrix}
\]
%
allora il vettore $X$ dovr\'a contenere $\begin{bmatrix}1 & 0 & 0\end{bmatrix}$
(poich\'e ciascun elemento della prima riga ha almeno una cifra pari, mentre
la seconda e la terza riga violano tale propriet\`a).
%
\begin{mdframed}[backgroundcolor=verylightgray]
  \begin{lstlisting}[language=myC]
#include <stdio.h>
#define N 3
#define M 4

int main() {
  int A[N][M];
  int X[N];
  return 0;
}
  \end{lstlisting}
\end{mdframed}

\vspace*{1ex}
\begin{center}{\Large Esercizio 2} ($11$ punti)\\
  (\textbf{si consegni} il file \texttt{coppie.c})
\end{center}

Si completi il seguente programma \texttt{coppie.c}
in modo che, inserita da tastiera
una serie di valori interi positivi,
conti e stampi a video quante coppie di numeri consecutivi della serie
soddisfano la seguente propriet\`a:
la somma delle cifre decimali del primo numero \`e uguale alla cifra
decimale pi\`u significativa del secondo numero. La fine della sequenza
viene indicata con un qualsiasi numero non positivo.
%
\begin{mdframed}[backgroundcolor=verylightgray]
  \begin{lstlisting}[language=myC]
#include <stdio.h>

int somma(int);
int left(int);

int main() {
  // continuare il codice
  return 0;
}

// restituisce la somma delle cifre decimali di un numero intero positivo
int somma(int n) {}

// RICORSIVA: restituisce la cifra decimale piu' significativa di un intero positivo
int left(int n) {}
  \end{lstlisting}
\end{mdframed}

In particolare, si completino le funzioni:
%
\begin{itemize}
\item \texttt{somma()}, che riceve un intero positivo $n$ e ritorna la somma di tutte
  le cifre decimali di $n$;
\item \texttt{left()} (\textbf{ricorsiva}), che riceve un intero positivo $n$
  e ritorna la sua cifra decimale pi\`u significativa.
\end{itemize}

Per esempio, se la sequenza inserita da tastiera fosse:
%
\begin{mdframed}[backgroundcolor=verylightgray]
\begin{verbatim}
12 23 53 245 53 12 0
\end{verbatim}
\end{mdframed}
%
allora il programma stamper\`a a video il valore $1$ perch\'e solo la coppia
$(23,53)$ di numeri consecutivi nella sequenza
soddisfa la propriet\`a $\texttt{somma}(23)=2+3=5=\texttt{left}(53)$.
%
\vspace*{1ex}
\begin{center}{\Large Esercizio 3} ($11$ punti)\\
  (\textbf{si consegni} il file \texttt{lista.c})
\end{center}

Si completi il seguente programma in modo che, data in input una sequenza di stringhe
di lunghezza non prefissata a priori (ciascuna di massimo $30$ caratteri)
inserita da tastiera e terminante con la stringa \texttt{END}
(che non fa parte della sequenza),
inizializzi una lista ottenuta trasformando ciascuna stringa
in un nodo che memorizza tre informazioni intere:
%
\begin{enumerate}
\item se la stringa \`e palindroma oppure no
  ($0$ = no, $1$ = s\`{\i});
\item quante cifre decimali pari contiene la stringa;
\item quante lettere maiuscole contiene la stringa.
\end{enumerate}
%
La sequenza termina quando l'utente inserisce
la stringa \texttt{END}.

\begin{mdframed}[backgroundcolor=verylightgray]
  \begin{lstlisting}[language=myC]
#include <stdio.h>
#include <stdlib.h>
#include <string.h>
#define END "END"

typedef struct nodo {
  int palindroma;
  int cifre;
  int maiuscole;
  struct nodo *next;
} node;

// prototipi, e' possibile aggiungere funzioni
void visualizza(node *);
int palindroma(char *);
int cifre(char *);
int maiuscole(char *);
node *inserisciincoda(node *, int, int, int);

int main() {
  node *l = NULL;
  char val[30+1];

  scanf("%s",val);
  while (strcmp(val, END) != 0) {
    int p = palindroma(val);
    int c = cifre(val);
    int m = maiuscole(val);
    l = inserisciincoda(l, p, c, m);
    scanf("%s", val);
  }
	
  visualizza(l);
  return 0;
}

void visualizza(node* lista) {
  while (lista != NULL){
    if (lista->next != NULL)
      printf("{%d, %d, %d} -> ", lista->palindroma, lista->cifre, lista->maiuscole);
    else
      printf("{%d, %d, %d}", lista->palindroma, lista->cifre, lista->maiuscole);
    lista = lista->next;
  }
  printf("\n");
}

node *inserisciincoda(node *l, int pal, int cif, int mai) {}
int palindroma(char *s) {}
int cifre(char *s) {}
int maiuscole(char* s) {}
  \end{lstlisting}
\end{mdframed}

Se, ad esempio, la sequenza di stringhe inserite fosse la seguente:
%
\begin{mdframed}[backgroundcolor=verylightgray]
\begin{verbatim}
casa Pesa int64 anna comPut3r END
\end{verbatim}
\end{mdframed}
%
allora il programma creer\`a e poi stamper\`a a video la seguente lista di cinque nodi:
%
\begin{mdframed}[backgroundcolor=verylightgray]
\begin{verbatim}
{0, 0, 0} -> {0, 0, 1} -> {0, 2, 0} -> {1, 0, 0} -> {0, 0, 1}
\end{verbatim}
\end{mdframed}

A tale scopo, si definiscano le funzioni:
%
\begin{itemize}
\item \texttt{palindroma()}, che riceve in ingresso una stringa e ritorna $1$
  se essa \`e palindroma (si considerino i caratteri maiuscoli diversi da quelli minuscoli)
  e $0$ altrimenti;
\item \texttt{cifre()}, che riceve in ingresso una stringa e ritorna il numero di
  cifre decimali pari in essa contenute;
\item \texttt{maiuscole()}, che riceve in ingresso una stringa e ritorna il numero di
  lettere maiuscole in essa contenute;
\item \texttt{inserisciincoda()}, che riceve in ingresso un puntatore \texttt{l}
  alla testa di una lista,
  tre informazioni intere da mettere nel nuovo nodo aggiunto in fondo ad \texttt{l},
  e restituisce il puntatore alla lista risultante dall'aggiunta del nuovo nodo
  in fondo ad \texttt{l}.
\end{itemize}

\end{document}
